% for Lua-based timestamping and basic timing
\RequirePackage{luacode}
% --- LUA BACKEND ---
\begin{luacode*}
-- Record CPU start time
compilation_start_time = os.clock()

-- Function for logging an elapsed event timestamp to the log file
function log_event(message)
  local current_time = os.clock()
  local elapsed = current_time - compilation_start_time
  local log_message = string.format("[+%.3fs] %s", elapsed, message or "")
  texio.write_nl("log", log_message)
end

-- Function for logging wall-clock time directly to the log file
function log_wall_clock_time(msg)
  local t = os.date("%Y-%m-%d %H:%M:%S")
  texio.write_nl("log", "[" .. t .. "] " .. (msg or ""))
end
\end{luacode*}

% --- LATEX FRONTEND ---
% 1. Log the start time immediately
\directlua{
  if log_wall_clock_time then
    log_wall_clock_time("Compilation started")
  end
}

% 2. Log both wall-clock and total CPU duration at the end of the run
\AtEndDocument{%
  \directlua{
    if log_wall_clock_time then
      log_wall_clock_time("Compilation finished")
    end
    if log_event then
      log_event("Total CPU time to end of document")
    end
  }%
}

% Command for explicit, manual logging
\newcommand{\ReportTimeStamp}[1]{\directlua{log_event("#1")}}

% Use the modern, built-in hook to log a message for each page
\AddToHook{shipout/before}{\ReportTimeStamp{Shipping out page \thepage}}
